\chapter{株式のロング・ショートとファクター}
\label{ch:equityfactor}

株式のファクター（\chapref{ch:factor}）は、クオンツのヘッジファンドの中心の戦略である。この章では、代表的なファクターの発見の論文、その説明、崩れた例、公表の後の減衰の研究をまとめる。

\section{代表的なファクターの発見}

\begin{table}[H]
\centering\small
\caption{株式の代表的なファクターと、その発見・検証の論文}
\label{tab:factorpapers}
\begin{tabular}{@{}P{0.16\linewidth}P{0.4\linewidth}P{0.36\linewidth}@{}}
\toprule
ファクター & 中身 & 主な論文 \\
\midrule
規模 & 時価総額の小さい株は、ベータで説明できないほど平均リターンが高い。 & Banz（1981）\cite{banz1981}、Fama・French（1993）\cite{ff1993} \\
割安 & 簿価・利益・キャッシュフローに比べて株価の低い株が、高い株を上回る。 & Lakonishok・Shleifer・Vishny（1994）\cite{lsv1994}、Fama・French（1993） \\
モメンタム & 過去 3〜12 か月の勝ち組が、その後 3〜12 か月も負け組を上回る。 & Jegadeesh・Titman（1993）\cite{jt1993}、Carhart（1997）\cite{carhart1997} \\
収益性・クオリティ & 収益性が高く、安定し、負債の少ない株が上回る。 & Novy-Marx（2013）\cite{novymarx2013}、Asness・Frazzini・Pedersen（2019）\cite{qmj2019} \\
投資 & 資産を大きく増やした会社は、その後のリターンが低い。 & Fama・French（2015）\cite{ff2015}、Hou・Xue・Zhang（2015）\cite{hxz2015} \\
低リスク & ベータや変動率の低い株が、リスクに比べて高いリターンを持つ。 & Frazzini・Pedersen（2014）\cite{bab2014}、Ang ほか（2006）\cite{ahxz2006} \\
会計の質 & 利益に占める発生主義の部分（アクルーアル）が大きい会社は、その後のリターンが低い。 & Sloan（1996）\cite{sloan1996} \\
\bottomrule
\end{tabular}
\end{table}

\subsection{モメンタム}

\begin{paperbox}[Jegadeesh・Titman（1993）\cite{jt1993}]
米国株（1965〜1989 年）で、過去 3〜12 か月のリターンで並べた勝ち組を買い、負け組を売る戦略が、その後 3〜12 か月の保有で有意な正のリターンを生むことを示した。過去 6 か月で選び 6 か月持つ戦略は、年率で平均 12\% 程度の超過リターンを上げた。この収益は市場のリスクでは説明できず、保有の 2 年目以降には一部が戻る。
\end{paperbox}

モメンタムは、その後、ほぼすべての国の株式と、債券・通貨・商品でも確認された（Asness・Moskowitz・Pedersen 2013 \cite{amp2013}、\ref{sec:amp} 節）。収益の説明は、主に行動の偏りによる。

\begin{itemize}
\item \textbf{反応の不足}：新しい情報が、少しずつしか価格に入らない。Hong・Stein（1999）\cite{hongstein1999} は、情報が投資家の間にゆっくり広がる市場では、価格が遅れて反応し、それを追う投資家が後で行き過ぎを生むことをモデルにした。
\item \textbf{注目の低い銘柄で強い}：Hong・Lim・Stein（2000）\cite{hls2000} は、規模を調整しても、アナリストの少ない銘柄ほどモメンタムが強く、特に悪い知らせ（負け組）の反映が遅いことを示した（「悪い知らせはゆっくり伝わる」）。
\end{itemize}

\begin{projbox}
Hong・Stein（1999）と Hong・Lim・Stein（2000）は、LLM-BLM の小型株の実験 S1 の仮説（注目の低い小型株では、決算発表文の情報の反映が遅れ、発表後のドリフトが残る）の直接の根拠である。
\end{projbox}

\subsection{割安}

Lakonishok・Shleifer・Vishny（1994）\cite{lsv1994} は、割安株（簿価やキャッシュフローに比べて株価の低い株）が割高株を上回るのは、割安株がリスクの高い株だからではなく、投資家が過去の業績の悪さを将来に引き伸ばしすぎ（外挿）、安く売りすぎるからだと論じた。一方、Fama・French は、割安株は苦境の会社のリスクの対価だと解釈する（\chapref{ch:factor}、表~\ref{tab:riskmisp}）。

\subsection{低リスク}

\begin{paperbox}[Frazzini・Pedersen（2014）「ベータに逆らって賭ける」\cite{bab2014}]
レバレッジを使えない（あるいは制限された）投資家は、高いリターンを求めて高ベータの資産を買う。その結果、高ベータの資産は割高に、低ベータの資産は割安になる。低ベータの資産をレバレッジで買い、高ベータの資産を売る BAB（betting against beta）のファクターが、米国株・国際株・国債・社債・通貨・商品のほとんどで正のリスク調整後のリターンを持つことを示した。また、資金の制約が厳しくなると BAB は損をする。
\end{paperbox}

Ang・Hodrick・Xing・Zhang（2006）\cite{ahxz2006} は、銘柄固有の変動率の高い株の平均リターンが著しく低いことを示した（固有のボラティリティのパズル）。

\subsection{収益性とクオリティ}

Novy-Marx（2013）\cite{novymarx2013} は、売上総利益を総資産で割った値が、簿価時価比率と同じくらい強く将来のリターンを予測し、しかも割安と負の相関を持つので、両方を組み合わせると成績が大きく良くなることを示した。Asness・Frazzini・Pedersen（2019）\cite{qmj2019} は、収益性・成長・安全性・配当などでクオリティを定義し、クオリティの高い株を買い低い株を売る QMJ のファクターが、多くの国で正のリスク調整後のリターンを持つことを示した。クオリティの高い株は、その良さに比べて十分に高い価格がついていない。

\section{あらゆる市場の割安とモメンタム}
\label{sec:amp}

\begin{paperbox}[Asness・Moskowitz・Pedersen（2013）\cite{amp2013}]
米国・英国・欧州・日本の株式、国別の株価指数、国債、通貨、商品の 8 つの市場で、割安とモメンタムの両方が正のリターンを持つことを示した。
\begin{itemize}
\item 割安とモメンタムは、どの市場の中でも負の相関を持つ。両方を組み合わせると、どちらか一方よりずっと安定する。
\item 市場をまたいで、割安どうし・モメンタムどうしは正の相関を持つ。共通の要因があり、その一部は資金の流動性のリスクに関係する。
\end{itemize}
\end{paperbox}

この論文は、ファクター投資が株式だけの現象ではなく、資産クラスをまたぐ一般的な現象であることを示し、多資産のファクター運用（オルタナティブ・リスク・プレミアムの商品）の理論的な根拠になった。

\section{崩れた例}

\subsection{モメンタムの急落}

\begin{paperbox}[Daniel・Moskowitz（2016）\cite{danielmoskowitz2016}]
モメンタムは、大きな下げ相場の後、市場が急に反発する局面で、数か月のうちに大きな損失を出す（1932 年、2009 年など）。下げ相場の後の負け組は、高ベータの株に偏っていて、反発の局面で急騰するからである。こうした急落は、事前の市場の状態と変動率からある程度予測でき、それに応じてモメンタムの量を動的に調整すると、シャープレシオが大きく改善することを示した。
\end{paperbox}

Barroso・Santa-Clara（2015）\cite{barroso2015} は、モメンタムの戦略を、その過去の実現の変動率で割って量を調整する（\chapref{ch:pc} のボラティリティ・ターゲット）だけで、急落がほぼ消え、シャープレシオがほぼ 2 倍になることを示した。

\subsection{割安の長い低迷と、クオンツの冬}

割安（HML）は、2007 年ごろから 2020 年まで長く損をし、特に 2018〜2020 年に大きく下げた。Blitz（2021）\cite{blitz2021} は、この時期を「クオンツの危機」と呼び、割安だけでなく、低ボラティリティ・小型・クオリティなど、多くのファクターが同時に損をしたこと、大型で割高な成長株への資金の集中（ビッグ・グロース）がその原因であったことを分析した。過去には割安が不調のときにモメンタムが補ってきたが、この時期には補いきれなかった、という指摘もある（Deep Research の答え 3 は、これを Blitz の報告として挙げた。原典の数値は確かめていない）。

\subsection{2007 年 8 月}

株式のマーケット・ニュートラルのクオンツが一斉に損をした 2007 年 8 月の出来事は、第 \ref{ch:statarb} 章と第 \ref{ch:failures} 章で扱う。

\section{公表の後の減衰と、再現性}

\subsection{公表の後に小さくなる}

\begin{paperbox}[McLean・Pontiff（2016）\cite{mcleanpontiff2016}]
学術誌に発表された 97 のリターンの予測の変数（アノマリー）について、論文の標本の期間の後、論文の発表の前（標本外）と、発表の後のリターンを比べた。
\begin{itemize}
\item 標本外では、標本の期間の中より平均 26\% 小さい。これは、統計の偏り（データ・マイニング）の上限の目安になる。
\item 発表の後は、標本の期間の中より平均 58\% 小さい。標本外との差（約 32\%）は、発表で情報を知った投資家の裁定による、と解釈される。
\item 標本の期間の中で強かったアノマリーほど、発表の後に大きく減る。裁定の費用が高い（小型、流動性が低い、固有の変動が大きい）銘柄のアノマリーほど、減り方が小さい。
\item 発表の後、アノマリーの銘柄の売買高や空売りが増える。
\end{itemize}
\end{paperbox}

\begin{figure}[H]
\centering
\begin{tikzpicture}[x=1mm,y=0.32mm]
\foreach \y in {0,25,50,75,100} {\draw[ink2!25] (0,\y)--(96,\y); \node[lab,left] at (0,\y) {\y};}
\fill[c1] (8,0) rectangle (26,100);   \node[lab,above] at (17,100) {100};
\fill[c1!60] (38,0) rectangle (56,74); \node[lab,above] at (47,74) {74};
\fill[c2!75] (68,0) rectangle (86,42); \node[lab,above] at (77,42) {42};
\node[lab,align=center,below] at (17,-2) {論文の標本の\\期間の中};
\node[lab,align=center,below] at (47,-2) {標本の後・\\発表の前};
\node[lab,align=center,below] at (77,-2) {発表の後};
\end{tikzpicture}
\caption{McLean・Pontiff（2016）の 97 のアノマリーの平均のリターンの減衰（論文の標本の期間の中を 100 とした値。要旨の 26\%・58\% から描いた）。標本の外での減少（26\%）は統計の偏りの上限の目安、発表の後の追加の減少は裁定による、と解釈される。}
\label{fig:mp}
\end{figure}

Jacobs・Müller（2020）\cite{jacobsmuller2020} は、39 か国の株式市場で 241 のアノマリーを調べ、公表の後の明確な減衰は米国だけで見られ、ほかの国では公表の後もリターンが残ることを示した。裁定の限界（空売りの難しさ、流動性、コスト）が、国による差を生んでいる可能性がある。Chordia・Subrahmanyam・Tong（2014）\cite{cst2014} は、米国で取引のコストが下がり売買が増えた 2000 年代以降（ティックの小数化の後）、アノマリーの多くが小さくなったことを示した。

\subsection{再現できるか}

ファクターの研究には、「そもそも本当にあったのか」という問いもある。

\begin{itemize}
\item \textbf{Hou・Xue・Zhang（2020）}\cite{hxz2020}：452 のアノマリーを、小型株の影響を抑えた同じ方法（NYSE の基準で分け、時価総額で加重）で作り直すと、6 割を超えるものが $t = 1.96$ の基準を超えなかった。多くのアノマリーは、極めて小さな株に依存していた。
\item \textbf{Jensen・Kelly・Pedersen（2023）}\cite{jkp2023}：153 のファクターを 93 か国で作り直し、ベイズの方法でファクターの間の情報を共有して評価すると、大多数のファクターが再現でき、米国の外でも成り立つ。「再現の危機」はないと結論した。
\item \textbf{Chen・Zimmermann（2022）}\cite{chenzimmermann2022}：300 を超える予測の変数のコードとデータを公開し、元の論文の結果の大部分が再現できることを示した。
\end{itemize}

二つの見方は、矛盾するというより、「何を再現とみなすか」の違いである。極めて小さな株を除き、時価総額で加重すると多くのアノマリーは弱くなるが、適切に作れば多くのファクターの方向と大きさは再現できる。

\subsection{取引コストの後}

Novy-Marx・Velikov（2016）\cite{novymarxvelikov2016} は、アノマリーの戦略の取引コストを推定し、回転率の低い戦略（割安、収益性など）はコストの後も残るが、回転率が月 50\% を超えるような戦略（短期の反転など）は、ほとんどがコストの後で有意でなくなることを示した。売買の頻度を下げる工夫（\chapref{ch:cost} の「売買しない領域」）で、コストを大きく減らせる。

Israel・Moskowitz（2013）\cite{israelmoskowitz2013} は、割安とモメンタムの収益が、空売りの側に依存しているわけではなく、大型株でも残り、時代とともに小さくなってもいない（規模を除く）ことを示した。運用会社の研究者による論文で、ファクター投資が現実に運用できることを論じたものである。

\section{確立していること、していないこと}

\begin{keybox}[まとめ]
\begin{description}[leftmargin=1em,style=nextline]
\item[確立している]
  割安・モメンタム・収益性・低リスクのファクターは、長い期間・多くの国・多くの資産クラスで正のリターンを持つ。アノマリーは、公表の後に平均して半分前後に小さくなる（米国）。極めて小さな株に依存するアノマリーは、実際には取れないことが多い。モメンタムはまれに大きく崩れ、変動率に応じた量の調整で和らげられる。
\item[議論が続いている]
  収益の原因がリスクか行動の偏りか。ファクターの収益がこれからも続くか（特に割安）。何百ものファクターのうち、独立な情報を持つものはいくつか（第 \ref{ch:ml} 章）。
\end{description}
\end{keybox}
